\section{Finite cuts with compatible ends}\label{sec:finite-cuts}

The stretching construction supplies separated sets and a scalar-curvature
reserve. We now describe the second input: a cutting procedure that retains
both sides of every cut and controls how later cuts meet them. Its output is
a family of smooth walls and complete manifolds obtained by attaching
cylinders to the closed sides. The definitions below make its hypotheses
and output explicit; Appendices~\ref{sec:packets}--\ref{sec:regularization}
give the proofs. Section~\ref{sec:countable} will separately justify its use
with countably many labels.

Let $X=M$ denote the original complete manifold, let $g$ be its original
metric, and write $d_0=d_g$ for its distance.  The metric used in an analytic
packet can change; the map to $X$ always records its position in the
original manifold.  Fix finitely many of the supplied ball pairs
\[
 A_i=\overline B_{d_0}(x_i,D),\qquad
 B_i=X\setminus B_{d_0}(x_i,2D),\qquad 1\le i\le\ell,
\]
with $A_i,B_i$ nonempty. There is one label for each ball pair; these
labels are also called \emph{colors}. An additional first label is radial.
Thus $N_*=\ell+1$, and the order is
radial, then $1,\ldots,\ell$.  Properness of $(X,g)$ makes each closed
outer ball compact.  No covering-space or quotient hypothesis is used.

The analytic input at stage zero is a complete smooth metric $h\ge g$ on $X$
and finitely many smooth positive circle warps for which the stabilized
scalar curvature is at least $.4$.  Each prescribed distance
$\dist_h(A_i,B_i)$ can be required to exceed any given finite number.
\cref{rounds:wide} supplies these data.  For finitely many requested
pairs its construction uses only finitely many compactly supported
increments and hence finitely many circle factors.  The countable
version has a locally finite list of nonunit warps and is used in
\cref{countable:construction}; the geometric statements below apply to
each finite collection of the relevant factors.  No uniform geometry
bounds on $h$, or uniform packetwise constants, are assumed.

For a single smooth cut, let $u$ be a smooth function on an
$m$-manifold $E$ without boundary, with $0$ a regular value, and put
$L=u^{-1}(0)$.
Attach $L\times[0,\infty)$ to each closed side $\{\pm u\ge0\}$
along its boundary, and also retain $L$ as a separate space.
The two extended sides have dimension $m$, while the wall has dimension
$m-1$. The cylinders replace the new boundaries by ends.
The following definition records these spaces and their incidences
after repeated cuts.

\begin{definition}[Packet diagram]\label{packet:category}
At stage $r$, a status is $\alpha\in\{+,-,*\}^{r}$.
Its entries record the first $r$ processed cuts.
A sign $+$ or $-$ selects the corresponding retained side of a cut,
and a star $*$ selects the wall, called a \emph{hit}.
The space $E_\alpha$ is a smooth second countable manifold without boundary,
possibly empty or disconnected; when nonempty its base dimension is
$m(\alpha)=n-\#\{i:\alpha_i=*\}$.  Its connected components are called
\emph{packets}.  Replacing signs by stars gives the face order.
If $B$ is a packet with $d$ additional stars relative to a coarser status,
then for that coarser status it has exactly one incident packet $E$ and
one corresponding occurrence of a full product tail
\begin{equation}\label{packet:tail}
 \iota_{BE}:B\times(T_1,\infty)\times\cdots\times(T_d,\infty)
       \longrightarrow E.
\end{equation}
The normal coordinates are labelled by the changed status entries.
Thresholds may depend on the incidence.  For a single entry, tails
from distinct face components are disjoint.  Intersections of sufficiently
far tails for different entries are exactly the tails over their common
further faces.  The identifications commute: composing face maps in
any order gives the same occurrence and the same individually labelled
normal coordinates.  After slightly raising the entry thresholds,
each closed tail is properly embedded.

Each packet $E$ is equipped with a continuous map $q_E:E\to X$.
After the radial stage we require the following additional properties.
Every packet has only finitely many incident faces and cofaces, including
through chains; the complement of its sufficiently far elementary tails
is compact.  On each sufficiently far face product these maps satisfy
$q_E(\iota_{BE}(b,\mathbf t))=q_B(b)$.
A metric, scalar function, warp, or set is \emph{compatible} if, in every
sufficiently far tail, it is exactly the product or pullback of the
corresponding face datum.  For metrics the normal factors are
$\dd t_1^2+\cdots+\dd t_d^2$; all warps are independent of the normal
coordinates there.  The thresholds for different data may differ.
\end{definition}

There is no boundary in an analytic packet.  Its face is represented by
an entire cylindrical end, not by a boundary subset.  Finite boundaries
appear only when we truncate these ends.

A \emph{regular split} applies the preceding construction using
compatible smooth defining functions whose zero sets are regular on
every status space. Each zero set is the new hit space.
In inherited ends the operation is the corresponding face operation
times the old normal factors. The proof that these operations preserve
the packet diagram is \cref{packet:split}.

For the diagrams constructed here, these maps are defined recursively,
starting with $q_X=\Id_X$.  In the split of \cref{packet:split}, let
$A$ be a new stage-$(r+1)$ packet with stage-$r$ parent $E$.
Define $p_A:A\to E$ by inclusion on a hit space or the bulk of a joined
side, and by $p_A(x,t)=x$ on each newly attached cylinder.
Set $q_A=q_E\circ p_A$.  Thus $q_A$ is pulled back from the base on
each fresh cylinder.  In a sufficiently far inherited tail of $A$,
let $C$ be the corresponding new face with old parent $B$, and let
$I$ denote the old normal intervals.  Under the identifications of
these tails with $C\times I$ and $B\times I$, the map $p_A$ is
$p_C\times\Id_I$; explicitly,
\[
 p_A\circ\iota_{CA}=\iota_{BE}\circ(p_C\times\Id_I),\qquad
 q_A\bigl(\iota_{CA}(c,\mathbf t)\bigr)=q_C(c).
\]
Thus the recursion preserves compatibility.  The proper maps $q_r$
on finite reconstructions are constructed in
\cref{packet:reconstruction}.

Figure~\ref{fig:regular-packet-split} depicts one extended side and the
separate hit space.
\begin{figure}[!b]
\centering
\begin{tikzpicture}[x=1cm,y=1cm,>=Stealth,font=\small]
  \fill[blue!7] (0,0) rectangle (2.3,2.2);
  \fill[orange!12] (2.3,0) rectangle (5,2.2);
  \draw[thick] (0,0) -- (5,0);
  \draw[thick] (0,2.2) -- (5,2.2);
  \draw[densely dashed,thick] (2.3,0) -- (2.3,2.2);
  \node[align=center] at (1.1,1.2) {retained\\closed side};
  \node[align=center] at (3.6,1.2) {$L\times[0,\infty)$\\attached\\cylinder};
  \draw[->,thick] (4.4,.5) -- (5.3,.5) node[right] {$s$};
  \node[align=center,below] at (2.3,-.1) {internal attachment seam};
  \draw[very thick,teal!70!black] (7,0) -- (7,2.2);
  \node[align=center,right] at (7.2,1.1) {separate hit space $L$\\base dimension $m-1$};
  \node[above] at (2.5,2.7) {either sign side: base dimension $m$};
\end{tikzpicture}
\caption{A regular cut retains each closed sign side with a full cylinder
attached, and also retains the wall as a separate hit space of one lower base
dimension. The attachment seam is internal to the extended side. In old
product ends this entire operation is a product. The raw metric and its
later smoothing are defined in Appendices~\ref{sec:cut} and~\ref{sec:regularization}.}
\label{fig:regular-packet-split}
\end{figure}


The label set and the status depth are finite in this section. The family
of connected packets may be countable; finiteness is required only for
the directed face and coface chains incident to each packet. The
downward diagram of a packet consists of that packet and all its further
faces. Numerical
requests are allocated downward through these chains before actual data
are constructed upward from faces. This does not require finite
components of the undirected incidence graph.

We now specify the analytic data that must survive each split.
For the finite-label statements in this section and the appendices we set $\sigma=1$
and use the following constants:
\begin{equation}\label{packet:constants}
\begin{gathered}
 \lambda_r=\frac14-\frac{r}{8N_*},\qquad
 b_r=1.12-\frac{.05r}{N_*},\qquad
 a_r=.60+\frac{.03r}{N_*}\quad(0\le r\le N_*),\\
 \pi_0=\frac{1}{10N_*},\qquad k=399,\qquad
 C=\frac{4k}{k+1}=3.99,\qquad
 e_0=\frac{\lambda_r-\lambda_{r+1}}{100}
     =\frac{1}{800N_*}\quad(r<N_*).
\end{gathered}
\end{equation}
Here $k$ is the number of circle factors added at each split; the initial
number may be different.  In particular $\lambda_r\ge1/8$, $a_r\le.63$,
and $1.07\le b_r\le1.12$.  The common floor $\pi_0$ and the equally
spaced stage parameters are bookkeeping for a finite list, not uniform
lower bounds for the countable construction.

\begin{definition}[Analytic invariant]\label{packet:invariant}
At stage $r$, write $J_r$ for the unprocessed labels.  On a packet $E$
of base dimension $m$, retain its base metric and all circle warps.
Write $R_E$ for the scalar curvature of that full warped metric and
$\omega_E$ for the product of its circle lengths, suppressing fixed
circle-volume constants.  Gradients of invariant functions are base
gradients; their integrals use $\omega_E\,\dd\vol_E$, and their
Laplacian is the corresponding weighted base Laplacian.
For every $j\in J_r$ the data comprise a compatible continuous profile
$\mu_j:E\to[-\infty,\infty]$ with locally Lipschitz finite truncations,
a compatible smooth potential $\Pi_j\ge\pi_0$, and compatible open
cores $U_j^+,U_j^-$ satisfying
\begin{equation}\label{packet:core-poles}
 \overline{U_j^\pm}\subset
 \operatorname{int}_E\{\mu_j=\pm\infty\}.
\end{equation}
On the open finite band the profile obeys, almost everywhere,
\begin{equation}\label{packet:profile}
 2|\grad\mu_j|\le a_r\mu_j^2+\Pi_j.
\end{equation}
Set
\begin{equation}\label{packet:required-set}
 G_j=E\setminus(U_j^+\cup U_j^-),\qquad
 D_E=\bigl\{x:\#\{j\in J_r:x\in G_j\}\ge m-1\bigr\}.
\end{equation}
The required scalar bound is
\begin{equation}\label{packet:curvature}
 R_E\ge\lambda_r+b_r\sum_{j\in J_r}\Pi_j
       \quad\hbox{on }D_E.
\end{equation}
For an unprocessed color $i$, the cores contain
$q_E^{-1}(A_i)\subset U_i^+$ and
$q_E^{-1}(B_i)\subset U_i^-$.  Enlarging cores must preserve these
conditions.  The sets $D_E$ are closed: they are finite unions of finite intersections
of the closed sets $G_j$ (and equal $E$ when the count threshold is
nonpositive).  All data are exact products
in the sense of \cref{packet:category}; at stage zero the single
space is $X$, with no product-end requirement.
\end{definition}

The threshold $m-1$ is preserved by the two kinds of output of a
split. A side keeps its dimension and loses one unprocessed label;
a hit loses both one dimension and that label. The lower threshold
$m-2$ defines a guard containing the projections of all new required
sets, even after the remaining cores have been enlarged.

\begin{lemma}[The two count inclusions]\label{packet:counts}
If $B$ is a $d$-fold further face of $E$, the product projection sends
$D_E$ in that tail into $D_B$.
At a split of label $j_0$ on a parent of dimension $m$, suppose the
remaining cores have been enlarged, giving $G'_j\subset G_j$ for
$j\ne j_0$, and define the closed guard
\begin{equation}\label{packet:guard}
 H_E^*=\bigl\{x:\#\{j\in J_r\setminus\{j_0\}:x\in G'_j\}
                      \ge m-2\bigr\}.
\end{equation}
Suppose a regular zero wall $L$ lies in $\operatorname{int}G_{j_0}$.
Pull the remaining cores to a new side or hit packet $A$ by bulk
inclusion, wall restriction, and projection of the attached cylinder.
Define $D_A$ using these cores, the remaining labels, and the dimension
of $A$, as in \eqref{packet:required-set}.
Then $D_A$ projects into both $H_E^*$ and the parent's old set $D_E$.
Guards likewise project into guards on further faces.
\end{lemma}

\begin{proof}
In an exact product tail all label memberships pull back unchanged,
whereas the face dimension is $m-d$.  A count at least $m-1$ is at
least $(m-d)-1$, proving the first assertion.  The same argument
with thresholds $m-2$ and $(m-d)-2$ proves the guard assertion.
On a new side, whose dimension is $m$, a point of $D_A$ has at least
$m-1$ remaining undecided labels.  This puts its projection in the
guard and, because $G'_j\subset G_j$, in the old required set.
On a hit packet, whose dimension is $m-1$, the required count is
$m-2$.  Its projection lies on $L\subset G_{j_0}$, so the processed
label contributes one further old undecided label.  The old count
is therefore at least $m-1$.  The cylinder projection has exactly
the same remaining memberships as its base, so the side argument
also covers all cylinder points.  These are pointwise count
statements and use no scalar estimate away from $D_E$.
\end{proof}

\begin{lemma}[Initialization and radial confinement]\label{packet:initialization}
Let the stage-zero stabilized scalar curvature be at least $.4$, and
suppose that every color pair has separation
\[
 \dist_h(A_i,B_i)>T_0+4,\qquad
 T_0=2\pi\sqrt{2/\pi_0}.
\]
Then the data of \cref{packet:invariant} exist at stage zero with
$\Pi_j=\pi_0$.  On a noncompact $X$ the radial cores can be chosen
as alternating separating subbands of a smooth proper function,
with compact smooth barrier domains around their closures lying in
the same pole plateaus.  Every regular radial partition respecting
these signs has locally finite compact wall components and compact
closed-side components before collar attachment.  On a compact $X$
a decided radial label with empty wall has the same properties.
\end{lemma}


\begin{theorem}[Finite compatible cutting]\label{finite:cut}
Fix a finite number $N_*$ of labels and the stage constants in
\eqref{packet:constants}. Suppose a stage-$r$ packet diagram is obtained by
the regular splitting construction of \cref{packet:split}, with compatible
smooth base metrics and a finite common list of positive circle warps,
and satisfies \cref{packet:invariant}. After the radial stage, require the
complete compact-part/product-end geometry of \cref{packet:category,packet:geometry}.
Let $j_0$ be the next label, and suppose every nonempty parent packet has
base dimension at least two. Then $j_0$ can be processed as follows.

On every parent there is a compatible smooth regular defining function
with the prescribed signs on the current forced core closures. Keep its
zero wall as the new hit space. To each closed sign side attach the full
product cylinder over its boundary. These pieces form a compatible smooth
stage-$(r+1)$ packet diagram with the continuous projections described
above. The wall on a parent $E$ lies in $\operatorname{int}G_{j_0}$.
The remaining profiles may first be modified, their cores enlarged,
and their potentials increased in the parents; these data are then carried to the new pieces by restriction and cylinder projection.
There are compatible smooth metrics and potentials, and exactly $399$
new positive circle factors on each new packet, for which the entire
stage-$(r+1)$ invariant holds. Earlier forced signs are preserved.
The new packets have complete compact-part/product-end geometry. This
assertion includes a possible one-dimensional output, on which its scalar
lower bound holds everywhere.

The operation fixes every old warp identically equal to $1$. For a
remaining label that is passive on the finite downward source diagrams
and the face copies used in their comparisons---constant pole profile, full corresponding core, and spatially
constant potential $c_j$ agreeing on compared faces---the construction
can preserve the pole and core and choose its new potential to be exactly
$(1+3\eta_r/2)c_j$, where $\eta_r=1/(200N_*)$.

At stage zero the single parent may instead be an arbitrary complete
smooth manifold with no product-end or bounded-geometry assumption. Subject to the same scalar and separation assumptions, if its
data are the radial initialization of \cref{packet:initialization}, the
same conclusion holds for the first split. Every resulting wall component and every closed-side component is compact before the cylinders are attached, and these components form locally finite
families.
\end{theorem}

\begin{proof}
Realize the processed profile by a smooth finite pressure and take a
local absolute minimizer of weighted perimeter minus its pressure-volume
term, as in \cref{cut:pressure,cut:bv-compactness}. The phase has the
forced signs of \cref{cut:barriers}; its boundary, called the true cut,
has the regular-part geometry of \cref{cut:gmt} but may still have
singularities. After the
radial cut, the generic pressure choice in
\cref{cut:generic-uniqueness} makes these true cuts converge to the
corresponding face products at packet ends
(\cref{cut:end-convergence}). On the initial radial parent use shift
zero and a phase from compactness; compact barriers replace this
post-radial growth and uniqueness argument.

The wall, and each retained side with its attached cylinder,
carry the quadratic-form lower bounds of \cref{cut:raw-domination}.
These bounds include a positive gradient reserve and, on a joined
side, the scalar contribution of its internal attachment seam.
They are not yet pointwise scalar bounds on smooth new packets.
Keep these true cuts and their pressures fixed. The constructions of
\Cref{local:anchors,local:sharp,shield:global} change only the remaining
profiles, enlarge their cores, and add nonnegative potentials. Their
form cost is paid by the gradient reserve and the gap
$\lambda_r-\lambda_{r+1}$, and the updated guard $H_E^*$ of
\eqref{packet:guard} misses the singular set. Thus the future required
regions can be controlled through regular portions of the true cut.
When the parent dimension is at most seven the
singular set is empty, and no singular forcing is needed; compactness on
the resulting smooth windows is the smooth case of
\cref{shield:compactness}. More generally this smooth-window case applies
whenever the limiting wall is smooth, including products of smooth faces
with Euclidean factors even when their total dimension is at least eight.
At the first radial split the same arguments
are applied on the locally finite compact transition blocks, as specified
in \cref{shield:global}.

The guarded approximation in \cref{regular:partition} now replaces
the true cut by a smooth regular wall. Its defining functions are
exact products in the old ends. The comparison with the true cut on
the guard, followed by the seam smoothing of \cref{regular:seam},
transfers the form bound to the smooth split. The positive
supersolutions of \cref{regular:correction,regular:solutions}
then supply $399$ compatible new circle factors; the scalar identity
in \cref{regular:step} converts that form bound into the pointwise
stage-$(r+1)$ scalar inequality. Throughout this transfer the parent
scalar inequality is used only over $D_E$: \cref{packet:counts}
places every new required region inside that old region as well as
inside the guard.

\Cref{packet:split,packet:geometry} give the
new geometry; \cref{packet:initialization} gives the radial confinement.
Unit entries are fixed by restriction, product extension, and seam
smoothing. For a remaining label passive in the sense stated above,
we omit its forcing updates, and
\cref{regular:constant-preservation,regular:profiles} reproduce its chosen
constant potential exactly. The choice of $\eta_r$ satisfies the strict
stage margins in \eqref{regular:eta}.
\end{proof}

This theorem is a finite-label statement. Its common positive floor and finite
list of warp factors do not provide a common floor or a global finite
list for the construction that follows. The next section proves the
stronger locality property actually needed there: on each whole packet
and each directed comparison diagram, only finitely many labels and warp
factors are nonconstant, and the remaining potential sum is a spatial
constant reproduced exactly by the construction.
